\hypertarget{appendix-5522ad176ca0}{}
\section*{Field reordering and AoS versus SoA}
\section{Field reordering and AoS versus SoA}\label{field-reordering-and-aos-versus-soa}

\textbf{Curriculum depth:} Introduced · \textbf{First seen:} C003 · \textbf{Latest development:} C003

\subsection{Definition and mechanism}\label{definition-and-mechanism}

Reordering fields can reduce alignment holes: conventional 1/8/1/4 members shrink 24→16 bytes when arranged 8/4/1/1. AoS (array of structures) interleaves full records; SoA (structure of arrays) separates each field into its own array. Smaller records do not prove faster code; hot fields and access patterns matter, and layout changes may break ABI.

\subsection{Prerequisites}\label{prerequisites}

\begin{itemize}
\tightlist
\item
  \href{}{Internal and tail padding (01-Foundations-I, appendix-f52f8c5a3afd)} --- The position and amount of unused layout space determine reordering/packing effects.
\item
  \href{}{Cache lines and memory bandwidth (01-Foundations-I, appendix-e4ccdbab8320)} --- Field adjacency and stride affect useful bytes fetched per access.
\end{itemize}

\subsection{Used by}\label{used-by}

Update downstream uses when later chunks build on this concept.

\subsection{Recall and next depth}\label{recall-and-next-depth}

Calculate footprint savings, then ask which fields a loop actually reads.

\textbf{Pending:} Measured AoS/SoA/AoSoA trade-offs, ECS and SIMD.

\subsection{Source and navigation}\label{source-and-navigation}

\href{Continuous\%20Notes/Computing\%20Foundations\%20-\%20From\%20Bits\%20to\%20C++\%20and\%20Systems\%20APIs\#Where\%20the\%20bytes\%20live}{C003 teaching chapter} · \href{}{Part definitions (01-Foundations-I, appendix-fd090d982a15 § M003 - Parts 41-60)} · \href{}{Dependency Map (02-Foundations-II, appendix-4ad47bdd269a)}
